\subsection{\texorpdfstring{格罗滕迪克群 \(R_{K}(A)\)}{格罗滕迪克群 R\_\{K\}(A)}}

设 \(\mathrm{K}\) 是一个交换域，\(\mathrm{A}\) 是一个 \(\mathrm{K}\)-代数。在 \(\mathrm{K}\) 上为有限维向量空间的 A-模的类所成之集是遗传的。相应的格罗滕迪克群记作 \(\mathrm{R}_{\mathrm{K}}(\mathrm{A})\)。它是一个自由 \(\mathbf{Z}\)-模，以族 \(([\mathrm{S}])_{\mathrm{S}\in \mathcal{S}}\) 为基，其中 \(\mathcal{S}\) 是在 \(\mathrm{K}\) 上为有限维的单 A-模的类所成之集。存在一个同态

\[
\dim : \mathrm{R} _ {\mathrm{K}} (\mathrm{A}) \longrightarrow \mathbf {Z}
\]

它由条件 \(\dim([E]) = [E : K]\)（对每个在 K 上为有限维的 A-模 E）刻画。当 A = K 时，它是一个同构。有效元所成的子幺半群记作 \(\mathrm{R}_{\mathrm{K}}(\mathrm{A})^{+}\)。

引理 2. —— 设 M 与 M' 是在 K 上为有限维向量空间的 A-模。M 与 M' 的支集（VIII, 第 190 页）不相交当且仅当存在 \(a \in A\)，使得 \(a_{M} = 0_{M}\) 且 \(a_{M'} = 1_{M'}\)。

假设存在 \(a \in \mathrm{A}\)，使得 \(a_{\mathrm{M}} = 0_{\mathrm{M}}\) 且 \(a_{\mathrm{M}'} = 1_{\mathrm{M}'}\)。设 S 是一个单 A-模。若 \(\operatorname{cl}(S)\) 属于 M 的支集 \(\mathcal{S}_{\mathrm{M}}\)，则 A-模 S 同构于 M 的某条若尔当–赫尔德列的商之一，且我们有 \(a_{\mathrm{S}} = 0_{\mathrm{S}}\)。同样，若 \(\operatorname{cl}(S)\) 属于 \(\mathcal{S}_{\mathrm{M}'}\)（即 \(\mathrm{M}'\) 的支集），则我们有 \(a_{\mathrm{S}} = 1_{\mathrm{S}}\)。由此得出 \(\mathcal{S}_{\mathrm{M}}\) 与 \(\mathcal{S}_{\mathrm{M}'}\) 不相交。

反之，假设集 \(S_{M}\) 与 \(S_{M'}\) 不相交。由于 M 与 \(M'\) 在 K 上是有限维的，它们是有限的。类属于 \(S_{M} \cup S_{M'}\) 的每个单 A-模 S 在 K 上是有限维的，从而更是在体 \(\operatorname{End}_{\mathrm{A}}(\mathrm{S})\) 上有限维。由命题 4 的推论 1（VIII, 第 83 页），存在元素 \(b \in A\)，使得有 \(b_{S} = 0_{S}\)（相应地，\(b_{S} = 1_{S}\)）对每个类属于 \(S_{M}\)（相应地，\(S_{M'}\)）的单 A-模 S 成立。设 \((\mathrm{M}_{i})_{0 \leqslant i \leqslant n}\) 是 M 的一条若尔当–赫尔德列。由前所述，我们有 \(bM_{i} \subset M_{i+1}\)（对 \(0 \leqslant i < n\)），因而 \((b^{n})_{\mathrm{M}} = 0_{\mathrm{M}}\)。使得 \(((1-b)^{m})_{\mathrm{M}^{\prime}}=0_{\mathrm{M}^{\prime}}\) 的自然数 m 的存在性用同样方式证明。令 \(\mathrm{P}(\mathrm{X})=1-(1-\mathrm{X}^{n})^{m}\) 且 \(a=\mathrm{P}(b)\)。多项式 \(\mathrm{P}(\mathrm{X})\) 是 \(X^{n}\) 的倍式，故我们有 \(a_{M}=0_{M}\)，而多项式 \(1-\mathrm{P}(\mathrm{X})\) 是 \((1-\mathrm{X})^{m}\) 的倍式，故我们有 \(a_{M^{\prime}}=1_{M^{\prime}}\)。证明至此完成。

设 \(\mathrm{L}\) 是 \(\mathrm{K}\) 的一个扩张。若 \(\mathrm{M}\) 是一个在 \(\mathrm{K}\) 上为有限维的 A-模，则 \(\mathrm{M}_{(\mathrm{L})}\) 是一个 \(\mathrm{A}_{(\mathrm{L})}\)-模，且在 \(\mathrm{L}\) 上为有限维。此外，对 A-模的每个正合列

\[
0 \longrightarrow \mathrm{M} ^ {\prime} \longrightarrow \mathrm{M} \longrightarrow \mathrm{M} ^ {\prime \prime} \longrightarrow 0
\]

（由 A-模组成），由它经标量扩张导出的 \(A_{(L)}\) -模的序列

\[
0 \longrightarrow \mathrm{M} _ {(\mathrm{L})} ^ {\prime} \longrightarrow \mathrm{M} _ {(\mathrm{L})} \longrightarrow \mathrm{M} _ {(\mathrm{L})} ^ {\prime \prime} \longrightarrow 0
\]

是正合的（II, §7, No.~7, 第 308 页，命题 14）。因此，存在唯一的环同态 \(u: \mathrm{R}_{\mathrm{K}}(\mathrm{A}) \to \mathrm{R}_{\mathrm{L}}(\mathrm{A}_{(\mathrm{L})})\)，使得对每个在 K 上为有限维的 A-模 M，有 \(u([\mathrm{M}]) = [\mathrm{M}_{(\mathrm{L})}]\)（No.~5）。

定理 1. —— 上面定义的同态 \(u: \mathrm{R}_{\mathrm{K}}(\mathrm{A}) \to \mathrm{R}_{\mathrm{L}}(\mathrm{A}_{(\mathrm{L})})\) 是单射。设 \(\xi\) 是 \(\mathrm{R}_{\mathrm{K}}(\mathrm{A})\) 的一个元素。则 \(\xi\) 是有效的当且仅当 \(u(\xi)\) 是有效的。

引理 3. —— 设 S 与 T 是两个在 K 上为有限维的互不同构的单 A-模。\(\mathrm{A}_{(\mathrm{L})}\)-模 \(\mathrm{S}_{(\mathrm{L})}\) 与 \(\mathrm{T}_{(\mathrm{L})}\) 的支集不相交。

由引理 2，存在元素 \(a \in \mathrm{A}\)，使得 \(a_{\mathrm{S}} = 0\) 且 \(a_{\mathrm{T}} = 1\)。元素 \(1 \otimes a\)（它属于 \(\mathrm{A}_{(\mathrm{L})}\)）在 \(\mathrm{S}_{(\mathrm{L})}\) 上作用为 0，在 \(\mathrm{T}_{(\mathrm{L})}\) 上作用为 1。由引理 2，\(\mathrm{A}_{(\mathrm{L})}\)-模 \(\mathrm{S}_{(\mathrm{L})}\) 与 \(\mathrm{T}_{(\mathrm{L})}\) 的支集不相交。

我们来证明定理 1。用 S 表示在 K 上为有限维的单 A-模的类所成之集。族 \(([S])_{S \in \mathcal{S}}\) 是 Z-模 \(R_{K}(A)\) 的一个基。设 \(S \in S\)。\(A_{(L)}\)-模 \(S_{(L)}\) 非零，故其支集非空。设 \(S'\) 是该支集的一个元素。由 VIII, 第 190 页的引理 1，存在同态 \(f_{S'} : R_L(A_{(L)}) \to Z\)，使得 \(f_{S'}([E]) = \ell_{S'}(E)\) 对每个在 L 上为有限维的 \(A_{(L)}\)-模 E 成立。按构造我们有 \(f_{S'}([S_{(L)}]) \neq 0\)，且由引理 3，有 \(f([T_{(L)}]) = 0\) 对每个 \(T \in S - \{S\}\) 成立。于是我们证明了：\(R_L(A_{(L)})\) 中形如 \([S_{(L)}]\)（\(S \in S\)）的元素在 Z 上线性无关。由此推出同态 u 是单射。

设 \(S \in \mathscr{S}\)，且设 \(S'\) 是 \(S_{(L)}\) 的支集的一个元素。对每个 \(\xi \in R_K(A)\)，\(\xi\) 的以下标 {[}S{]} 为指标的坐标（在基 \(([S])_{S \in \mathscr{S}}\) 中）是 \(f_{S'}(u(\xi)) / [S_{(L)} : S']\)。由此得出：若 \(u(\xi)\) 是有效的，则 \(\xi\) 也是有效的。

